\subsection{The frozen CAR covariance}

The finite collar used by \texttt{v995} has eight complex fermion
modes, hence
\[
 \CAR_8\simeq M_{2^8}(\C)=M_{256}(\C).
\]
In dimensionless modular units the one-particle Hamiltonian is
\[
 h_{\mathrm{collar}}
 =\Delta\,\operatorname{diag}(1,2,\ldots,8),
\qquad
 \Delta=6\log(3/2),
\tag{3.6}\label{eq:collar-h}
\]
and the script sets $\beta=1$.  Therefore
\[
 q=e^{-\Delta}=(2/3)^6=\frac{64}{729},
\qquad
 C_j=\frac{q^j}{1+q^j}
 =\frac1{1+e^{j\Delta}},\quad j=1,\ldots,8.
\tag{3.7}\label{eq:frozen-C}
\]
This is exactly Araki's covariance
\eqref{eq:araki-cov}.  All eigenvalues lie strictly between $0$ and
$1$, so the finite CAR representation is factorial and the density
matrix is faithful.  On the occupation basis
$n=(n_1,\ldots,n_8)\in\{0,1\}^8$,
\[
 p(n)=Z^{-1}q^{\sum_{j=1}^{8}jn_j},
\qquad
 Z=\prod_{j=1}^{8}(1+q^j).
\tag{3.8}\label{eq:gibbs-weights}
\]
Thus all $256$ weights are positive, and
\[
 \frac{p(n+e_j)}{p(n)}=q^j
\]
whenever $n_j=0$.  These are the exact detailed-balance relations
checked in \texttt{v995}.  The angular convention
$\beta_{\mathrm{angle}}=2\pi$ is recovered by absorbing $2\pi$ into
the modular Hamiltonian: $\beta h=(2\pi)H_{\mathrm{angle}}$.

\begin{remark}[What Araki proves here]
Araki's theorem proves uniqueness of the KMS state of the full CAR
dynamics and its determination by \eqref{eq:frozen-C}.  Restricting
that state to the even/bilinear $\mathrm{SO}(16)_1$ subnet selects the
vacuum-sector thermal state.  It does \emph{not} say that the
observable $\mathrm{SO}(16)_1$ net has only one rotational KMS state:
that net has four DHR sectors and hence four extremal sector Gibbs
states.  Holomorphic extension to $\B$ is what removes this sector
choice.
\end{remark}

\subsection{Holomorphy and thermal classification}

By \eqref{eq:holomorphic} and Theorem~\ref{thm:klm},
\[
 1=\mu_{\B}
 =\sum_{\rho\in\operatorname{Irr}(\B)}d(\rho)^2.
\tag{3.9}\label{eq:one-sector}
\]
The vacuum already contributes one, so there can be no other
irreducible DHR sector.  Theorem~\ref{thm:lt} then has one Gibbs
state and no other extreme point.  This verifies the analytic
hypothesis needed for unconditional rotational KMS uniqueness.

\section{The holomorphic KMS extension theorem}

\begin{theorem}[Holomorphic rotational KMS extension]
\label{thm:main}
Let $\A\subset\B$ satisfy the hypotheses of Section~1: for every
interval $I$, $\A(I)\subset\B(I)$ is an irreducible index-four
inclusion, $\A=\B^{\Z_4}$ for a proper internal gauge action
commuting with rotations, and $\B$ is a holomorphic conformal net
with $\mu_{\B}=1$ and $c=8$.  Let $\beta>0$ and let
$\omega_\beta$ be the restriction \eqref{eq:omega-def} of the
vacuum-sector rotational Gibbs state.  Then:
\begin{enumerate}[label=\textup{(\alph*)}]
\item $\omega_\beta$ has a rotational $\beta$-KMS extension to
  $\B$.
\item The maps
  \[
    E_I=\frac14\sum_{r=0}^{3}\alpha_r:
    \B(I)\to\A(I)
  \]
  are the unique normal faithful conditional expectations.  They
  form a net expectation $E$, and
  \[
    \widetilde\omega_\beta
    =\omega_\beta\circ E
    =\Phi_\beta^{\B}.
  \tag{4.1}\label{eq:main-extension}
  \]
  This is the unique $\Z_4$-invariant state extension of
  $\omega_\beta$, hence in particular the unique
  $\Z_4$-invariant rotational KMS extension.
\item Every rotational $\beta$-KMS extension of $\omega_\beta$ is
  $\Z_4$-invariant.  Consequently
  $\widetilde\omega_\beta$ is the unique rotational KMS extension
  without an invariance qualifier.
\end{enumerate}
\end{theorem}

\begin{proof}
\emph{Part (a).}
The Gibbs functional \eqref{eq:gibbs} is a rotational
$\beta$-KMS state on $C^*_{\mathrm{ln}}(\B)$ by
Theorem~\ref{thm:lt}, or directly by the cyclicity of the trace on
analytic elements:
\[
\begin{aligned}
 \Phi_\beta^{\B}(x\tau_{t+i\beta}(y))
 &=Z_\beta^{-1}
   \Tr(e^{-\beta L_0}x e^{itL_0}e^{-\beta L_0}
        y e^{\beta L_0}e^{-itL_0})\\
 &=Z_\beta^{-1}
   \Tr(e^{-\beta L_0}\tau_t(y)x)
 =\Phi_\beta^{\B}(\tau_t(y)x).
\end{aligned}
\tag{4.2}\label{eq:gibbs-kms}
\]
Its restriction to a $\tau$-invariant $C^*$-subalgebra remains KMS:
the same strip-analytic functions work for $x,y\in
C^*_{\mathrm{ln}}(\A)$.  By definition that restriction is
$\omega_\beta$.  Hence $\Phi_\beta^{\B}$ is the required extension.

\emph{Part (b): expectation and its uniqueness.}
Section~3.1 proves directly that the group average $E_I$ is a normal
faithful conditional expectation.  If $F_I:\B(I)\to\A(I)$ is any
other normal faithful conditional expectation, the
Radon--Nikodym part of Theorem~\ref{thm:qsystem} gives
\[
 F_I(x)=E_I(h^{1/2}xh^{1/2})
\]
with $h\in\A(I)'\cap\B(I)$.  Irreducibility gives
$h=\lambda\id$.  Since $F_I(\id)=\id=E_I(\id)$, $\lambda=1$, and
$F_I=E_I$.  This also shows why the bimodule property in the
definition of a conditional expectation is load-bearing: without
it, the Radon--Nikodym theorem does not apply, exactly as the
finite mutants in \texttt{v993} demonstrate.

The Gibbs state is gauge invariant.  Indeed the internal symmetry
commutes with $L_0$ and is unitarily implemented, so
\[
\begin{aligned}
 \Phi_\beta^{\B}(\alpha_r(x))
 &=Z_\beta^{-1}\Tr(e^{-\beta L_0}V_rxV_r^*)\\
 &=Z_\beta^{-1}\Tr(V_r^*e^{-\beta L_0}V_rx)
 =\Phi_\beta^{\B}(x).
\end{aligned}
\tag{4.3}\label{eq:gauge-invariance}
\]
Consequently
\[
 \Phi_\beta^{\B}(E(x))
 =\frac14\sum_{r=0}^{3}\Phi_\beta^{\B}(\alpha_r(x))
 =\Phi_\beta^{\B}(x).
\tag{4.4}\label{eq:E-preserves}
\]
Because $E(x)\in\A$ and
$\Phi_\beta^{\B}|_{\A}=\omega_\beta$,
\[
 \Phi_\beta^{\B}(x)
 =\Phi_\beta^{\B}(E(x))
 =\omega_\beta(E(x)),
\]
which proves \eqref{eq:main-extension}.

Now let $\psi$ be any $\Z_4$-invariant state on $\B$ whose restriction
to $\A$ is $\omega_\beta$.  No KMS hypothesis is needed:
\[
 \psi(x)
 =\frac14\sum_{r=0}^{3}\psi(\alpha_r(x))
 =\psi(E(x))
 =\omega_\beta(E(x)).
\tag{4.5}\label{eq:inv-unique}
\]
Thus $\psi=\omega_\beta\circ E$.

\emph{The modular-theory chain.}
We spell out the Takesaki step because it is frequently reversed
incorrectly.  In the GNS representation of the faithful KMS state,
the modular group and the physical dynamics are related, with the
present KMS convention, by
\[
 \sigma_s^{\Phi_\beta^{\B}}(\pi_\Phi(x))
 =\pi_\Phi(\tau_{-\beta s}(x)).
\tag{4.6}\label{eq:kms-modular}
\]
Rotations preserve the subnet, hence
\[
 \sigma_s^{\Phi_\beta^{\B}}
   (\pi_\Phi(\A(I)))=\pi_\Phi(\A(I)).
\]
Takesaki's theorem now gives the unique
$\Phi_\beta^{\B}$-preserving normal faithful expectation.  By
\eqref{eq:E-preserves}, the group average is such an expectation, so
it is the Takesaki expectation.  Moreover,
\[
 \sigma_s^{\omega_\beta}
 =\sigma_s^{\Phi_\beta^{\B}}|_{\A(I)}
 =\tau_{-\beta s}|_{\A(I)},\qquad
 E_I\sigma_s^{\Phi_\beta^{\B}}
 =\sigma_s^{\omega_\beta}E_I.
\tag{4.7}\label{eq:modular-chain}
\]
Notice the logical order: existence of the Gibbs extension gives a
state on $\B$; modular invariance then identifies its
state-preserving expectation.  Takesaki's criterion alone would not
manufacture a KMS extension from an arbitrary state on $\A$.

\emph{Part (c): removal of the invariance qualifier.}
Let $\psi$ be a rotational $\beta$-KMS state on
$C^*_{\mathrm{ln}}(\B)$ extending $\omega_\beta$.  By
Theorem~\ref{thm:klm} and \eqref{eq:one-sector}, $\B$ has only the
vacuum irreducible DHR sector.  Theorem~\ref{thm:lt} says that every
rotational KMS state is a convex combination of Gibbs states in
irreducible sectors.  There is only one such state,
$\Phi_\beta^{\B}$, so
\[
 \psi=\Phi_\beta^{\B}.
\]
Equation \eqref{eq:gauge-invariance} makes this state
$\Z_4$-invariant, and part (b) then gives
$\psi=\omega_\beta\circ E$.  This proves unconditional uniqueness.
\end{proof}

\section{Where the mixing directions went}

\subsection{Relative-commutant mixing}

Fix an interval and write $N=\A(I)$, $M=\B(I)$.  An internal
multiplicity mixer for the inclusion is an operator $T\in N'\cap M$.
By irreducibility,
\[
 N'\cap M=\C\id,
\tag{5.1}\label{eq:relative-commutant}
\]
so $T=\lambda\id$.  A normalized perturbation cannot use this scalar
to create a new state direction.  Explicitly, if
$T=T^*$ and
\[
 \psi_\varepsilon(x)
 =\frac{\Phi_\beta^{\B}
 ((\id+\varepsilon T)x(\id+\varepsilon T))}
 {\Phi_\beta^{\B}((\id+\varepsilon T)^2)},
\]
then \eqref{eq:relative-commutant} makes numerator and denominator
carry the same scalar factor, and
\[
 \psi_\varepsilon=\Phi_\beta^{\B}.
\tag{5.2}\label{eq:scalar-pert}
\]
Thus every mixing direction that would be implemented by an
operator commuting with the subnet is trivial.

\subsection{Superselection-sector mixing}

A different possible mixing is a convex combination of thermal
states in inequivalent irreducible representations.  KLM gives the
sector count through
\[
 \mu_{\B}=\sum_jd(\rho_j)^2.
\]
Holomorphy makes the sum equal one, so only $\rho_0=\id$ remains.
Longo--Tanimoto identifies the extreme rotational KMS states with
the sector Gibbs states.  Therefore there is no second thermal
sector into which a state can mix.

These are two distinct exclusions:
\begin{center}
\setlength{\fboxsep}{7pt}
\fbox{\begin{minipage}{0.86\textwidth}
\noindent
$\A(I)'\cap\B(I)=\C\id$ removes local multiplicity mixing.

\smallskip\noindent
$\mu_{\B}=1$, together with the Longo--Tanimoto classification,
removes thermal-sector mixing.
\end{minipage}}
\end{center}
Neither exclusion should be substituted for the other.  In
particular, holomorphy alone is a statement about DHR sectors, not
an abstract theorem about all possible dynamics.

\section{TFPT corollary and the P1 reading}

\subsection{The finite matrix shadow}

\begin{corollary}[Interpretation of \texttt{v995}]
\label{cor:v995}
Assume that the TFPT collar scaling limit realizes the net inclusion
of Theorem~\ref{thm:main} with the stated rotation dynamics.  Then
the exact finite results of \texttt{v995} are finite-dimensional
shadows of the two uniqueness mechanisms:
\begin{enumerate}[label=\textup{(\roman*)}]
\item on $\CAR_8=M_{256}$, the KMS equations have rank
  $256^2-1=65535$ on the state tangent and kernel zero;
\item all $33$ selected same-$\Z_4$-charge, unequal-modular-energy
  Hermitian directions fail stationarity and are killed by KMS;
\item on the finite crossed product, fixing the grade-zero
  restriction leaves
  \[
    3\cdot256^2=196608
  \]
  real linear extension coordinates, and invariance under
  $E=(1/4)\sum_r\alpha_r$ kills all of them;
\item covariance is redundant after the full finite-factor KMS
  condition has been imposed.
\end{enumerate}
\end{corollary}

\begin{proof}
Items (i)--(iv) are exactly the finite linear-algebra statements of
\texttt{v995}.  The continuum interpretation is conditional on the
scaling-limit identification because finite matrix rank does not by
itself prove factoriality or classify locally normal KMS states of a
type-III net.  Under that identification, item (iii) is the finite
Fourier version of \eqref{eq:inv-unique}, while items (i), (ii), and
(iv) are the finite CAR instance of Theorem~\ref{thm:araki}.
\end{proof}

\subsection{Equidistribution over four deck channels}

The exact statement needed for the P1 reading is elementary but has
one hypothesis that must be visible.

\begin{lemma}[Transitive four-channel equidistribution]
\label{lem:equidistribution}
Let $p_0,p_1,p_2,p_3$ be orthogonal projections in a response
algebra with $\sum_rp_r=\id$, and suppose a $\Z_4$ action is
transitive:
\[
 \alpha_s(p_r)=p_{r+s\ {\rm mod}\ 4}.
\tag{6.1}\label{eq:transitive}
\]
If $\psi$ is $\Z_4$-invariant, then
\[
 \psi(p_0)=\psi(p_1)=\psi(p_2)=\psi(p_3)=\frac14.
\tag{6.2}\label{eq:quarter}
\]
If a normalized total modular response is
$R_{\mathrm{tot}}=1/\beta_{\mathrm{angle}}$ and is additive over
these channels, then each channel carries
\[
 R_r=\frac1{4\beta_{\mathrm{angle}}}.
\tag{6.3}\label{eq:response-quarter}
\]
\end{lemma}

\begin{proof}
For any $r,s$, invariance and transitivity give
\[
 \psi(p_{r+s})=\psi(\alpha_s(p_r))=\psi(p_r).
\]
All four values are therefore a common number $w$.  Normalization
gives $1=\psi(\id)=\sum_r\psi(p_r)=4w$, hence $w=1/4$.
Additivity of the response then gives
$R_r=R_{\mathrm{tot}}/4$.
\end{proof}

\begin{corollary}[The theorem-conditional P1 identity]
\label{cor:p1}
Under Theorem~\ref{thm:main} and the response-resolution hypothesis
of Lemma~\ref{lem:equidistribution}, the canonical geometric
normalization
\[
 \beta_{\mathrm{angle}}=2\pi
\tag{6.4}\label{eq:beta-angle}
\]
gives
\[
 c_3=R_r
 =\frac1{4\beta_{\mathrm{angle}}}
 =\frac1{4\cdot2\pi}
 =\frac1{8\pi},
\qquad
 4\beta_{\mathrm{angle}}c_3=1.
\tag{6.5}\label{eq:p1}
\]
\end{corollary}

\begin{proof}
Substitute \eqref{eq:beta-angle} into
\eqref{eq:response-quarter}.  No further analytic input is used.
\end{proof}

\begin{remark}[Exactly what P1 gains]
Theorem~\ref{thm:main} proves uniqueness and gauge invariance of the
rotational KMS extension for the abstract holomorphic net.  Once a
four-channel response resolution satisfying
\eqref{eq:transitive} is identified, equality of its weights is no
longer an assumption: it is Theorem~\ref{thm:main} plus
Lemma~\ref{lem:equidistribution}.  The factor $1/4$ is therefore a
theorem-conditional consequence of index four and transitivity.

What is \emph{not} proved is that the TFPT determinant-line response
is represented by projections $p_r$ of this kind, or that its
physical coefficient is the channel response $R_r$.  An outer
fixed-point inclusion of factors does not supply four central local
sector projections: the grading subspaces are not central summands.
That operator identification is the joint
\texttt{P1.INDEX.KMS.01} and
\texttt{ALPHA.QUILLEN.EXACT.01} bridge.
The angular normalization $\beta_{\mathrm{angle}}=2\pi$ is the
geometric Bisognano--Wichmann/modular input, not a consequence of
the group order.  Accordingly, \texttt{AX.P1.01} remains an axiom;
\eqref{eq:p1} is a theorem-conditional identity, not an unconditional
derivation of P1.
\end{remark}

